\section{Reconstruction-Then-Ridge Estimation of the Learned-Lifting Baseline}
\label{app:estimator}

\textbf{Why two stages, and why the intercept escapes the penalty.} A learned lifting is the
alternative a reader will reach for, so it has to be fitted well enough that a null result is
about the data. Fitting the encoder and the dynamics together lets one absorb the other's
error: a lifting free to move can make almost any dynamics look adequate, which turns the
comparison into a statement about optimization. Fitting the lifting first and solving the
dynamics in closed form afterwards removes that freedom. The intercept escapes the ridge
penalty for a separate reason: a penalized intercept is pulled toward zero, the dynamics block
absorbs the latent mean to compensate, and the estimator is no longer the one we describe.
\uline{Reconstruction-then-ridge therefore separates the two fits so that the learned lifting
cannot rescue a weak dynamics estimate, and leaves the intercept unpenalized so that the
dynamics block is not shrunk into carrying the mean.}

\textbf{Lift and solve.} Equation~\eqref{eq:ae-koopman} gives the lift and the latent map,
where $\xi_t$ is the lifted state and $K$, $B$, $c$ are its transition, control and intercept.
Equation~\eqref{eq:stage-two-solve} gives the second stage, where $\Lambda$ is the ridge
penalty and $\lambda_{\mathrm{rec}}$ weights the reconstruction term of the first.

\textbf{Two stages in order.} Stage one descends the reconstruction loss alone, leaving the
latent dynamics untouched. Stage two freezes the encoder, pushes every training pair through it
once, and solves for the latent dynamics in closed form. The ridge penalty carries a zero in the
position of the appended constant column, which is the unpenalized intercept; no comment in the
released code marks that zero, and a reader assuming a uniform penalty is reading a different
estimator. Appendix~\ref{app:joint} states the joint objective that descends reconstruction,
one-step and multi-step terms together. We report no ablation of its multi-step term, which
takes one gradient step per epoch against one per minibatch for the others, so its weight
confounds with dataset size and was never swept at comparable strength.
